\section{Вычислительный эксперимент. Принципы проведения вычислительного эксперимента. Модель, алгоритм, программа}

\subsection{Понятие вычислительного эксперимента}

\textbf{Вычислительный эксперимент} --- исследование математической модели объекта с помощью численных методов и программных средств. В отличие от натурного эксперимента непосредственно исследуется не реальный объект, а его модель.

Пусть модель имеет вид
\[
\mathcal M(u,p)=0.
\]
После дискретизации получаем конечномерную задачу
\[
\mathcal M_h(u_h,p)=0,
\]
где $h$ обозначает параметры сетки или иной аппроксимации. Серия расчётов при различных $p$ и $h$ позволяет изучать закономерности, чувствительность, критические режимы и прогнозы.

\subsection{Триада ``модель--алгоритм--программа''}

Основой вычислительного эксперимента является
\[
\boxed{
\text{математическая модель}
\longrightarrow
\text{вычислительный алгоритм}
\longrightarrow
\text{программа}}.
\]

\textbf{Модель} определяет уравнения, ограничения, параметры, начальные и граничные условия. Ошибка на этом уровне означает, что выбрано неадекватное описание объекта.

\textbf{Алгоритм} переводит математическую постановку в конечную последовательность вычислений. Для него важны аппроксимация, устойчивость, сходимость, вычислительная сложность и критерии остановки.

\textbf{Программа} реализует алгоритм на конкретной вычислительной системе. Здесь добавляются вопросы представления данных, программных ошибок, библиотек, параллельного исполнения и ввода-вывода.

Корректная программа не исправляет плохую модель, а хороший численный метод не исправляет программную ошибку; поэтому уровни проверяются раздельно.

\subsection{Основные этапы}

Типичный цикл:
\begin{enumerate}
\item постановка цели исследования;
\item построение математической модели;
\item исследование её корректности и характерных масштабов;
\item выбор и анализ численного метода;
\item программная реализация;
\item \textbf{верификация} алгоритма и программы;
\item планирование и проведение серии расчётов;
\item \textbf{валидация} модели по экспериментальным данным;
\item анализ результата, чувствительности и неопределённости;
\item при необходимости --- уточнение модели и повтор цикла.
\end{enumerate}

Таким образом,
\[
\text{модель}\rightarrow\text{расчёт}\rightarrow\text{сравнение}\rightarrow\text{уточнение}
\]
--- итерационный процесс.

\subsection{Планирование серии расчётов}

Если модель зависит от параметров
\[
p=(p_1,\ldots,p_m),
\]
нужно выбрать набор расчётных точек $p^{(1)},\ldots,p^{(N)}$. Простейший подход --- менять один фактор при фиксированных остальных; при большом числе факторов применяются методы планирования эксперимента, позволяющие извлечь больше информации из ограниченного числа запусков.

Локальная чувствительность выходной величины $y$ к параметру $p_i$ характеризуется
\[
S_i=\frac{\partial y}{\partial p_i}
\approx
\frac{y(p_i+\Delta p_i)-y(p_i)}{\Delta p_i}.
\]
План должен включать не только физические параметры, но и численные параметры: шаг сетки, временной шаг, допуски решателей.

\subsection{Верификация и сеточная сходимость}

\textbf{Верификация} отвечает на вопрос: ``правильно ли решаются уравнения модели?'' Основные средства:
\begin{itemize}
\item сравнение с аналитическим или изготовленным решением;
\item независимая реализация или другой численный метод;
\item контроль законов сохранения;
\item модульные и регрессионные тесты;
\item исследование сходимости при измельчении дискретизации.
\end{itemize}

Если метод имеет порядок $p$, ожидается
\[
\|u-u_h\|\approx Ch^p,
\]
и при переходе $h\to h/2$ ошибка асимптотически уменьшается приблизительно в $2^p$ раз.

Общая численная ошибка включает по меньшей мере погрешность дискретизации, незавершённости итераций и округления; кроме неё существует ошибка исходной математической модели и данных.

\subsection{Валидация и анализ результата}

\textbf{Валидация} отвечает на вопрос: ``адекватно ли выбранная модель описывает реальный объект?'' Для неё нужны экспериментальные/натурные данные. Например, сравнивают
\[
r_i=y_i^{\rm exp}-y_i^{\rm calc}
\]
и исследуют не только норму ошибки, но и систематическую структуру остатков.

Верификация должна логически предшествовать валидации: нельзя судить о физической модели по результату, если неизвестно, корректно ли решается её математическая постановка.

После серии расчётов анализируют зависимости от параметров, пространственно-временные поля, экстремумы, критические режимы и область применимости модели. Результат должен сопровождаться параметрами модели, сеткой, методом, допусками и оценкой точности.

\subsection{Воспроизводимость}

Для повторения эксперимента необходимо сохранить исходные данные, параметры, модель, версию программы и библиотек, настройки дискретизации и решателя, а для стохастических расчётов --- параметры генератора случайных чисел. Желательны контроль версий, сценарии автоматического запуска и сохранение метаданных.


\subsection{План вычислительного эксперимента}

Перед серией расчётов полезно сформулировать экспериментальный план:
\begin{enumerate}
    \item какие факторы $p_1,\ldots,p_m$ изменяются;
    \item в каких диапазонах они физически допустимы;
    \item какие выходы $y_1,\ldots,y_k$ измеряются;
    \item какой критерий сравнивает варианты;
    \item какая точность нужна для принятия решения.
\end{enumerate}

Изменение одного фактора при фиксированных остальных удобно для первичного
анализа, но не выявляет взаимодействий. Факторные планы и выборки в
многомерном пространстве параметров позволяют строить приближённый отклик
\[
y\approx \beta_0+\sum_i\beta_ip_i+\sum_{i<j}\beta_{ij}p_ip_j+\cdots
\]
и оценивать совместное влияние факторов.

\subsection{Разделение ошибок}

Вычислительный результат содержит по крайней мере три разных источника
неопределённости:
\[
\text{модельная} + \text{дискретизационная} + \text{алгоритмическая/округление}.
\]
Поэтому верификация программы не заменяет валидацию модели. Совпадение
результата с экспериментом также не доказывает, что каждая внутренняя часть
алгоритма реализована правильно: возможна компенсация ошибок.

Для сеточной сходимости удобно рассматривать последовательность $Q_h$ и
проверять асимптотический режим
\[
Q_h=Q+C h^p+o(h^p).
\]
Тогда значения на нескольких сетках позволяют оценивать наблюдаемый порядок и
экстраполировать результат.

\subsection{Чувствительность и неопределённость}

Если выход $y=F(p)$ зависит от параметров, локальная чувствительность задаётся
\[
S_i=\frac{\partial F}{\partial p_i}.
\]
При случайных параметрах с ковариационной матрицей $C_p$ линейная оценка
распространения неопределённости имеет вид
\[
C_y\approx J C_pJ^T,
\qquad J=\frac{\partial F}{\partial p}.
\]
Для существенно нелинейных моделей используют выборочные методы, например
Монте--Карло. В отчёте важно отличать численную погрешность конкретного расчёта
от неопределённости входных данных и параметров.

\subsection{Протокол воспроизводимого расчёта}

Научный вычислительный эксперимент должен быть повторяемым. Следует сохранять
версию модели и программы, исходные данные, параметры, сетку, критерии остановки,
версии библиотек, вычислительную среду и seeds генераторов случайных чисел.
Системы контроля версий и автоматизированные сценарии запуска превращают
единичный расчёт в проверяемую процедуру.


\subsection{Что важно сказать на экзамене}

Главное --- не сводить вычислительный эксперимент к ``запуску программы''. Нужно показать полный цикл
\[
\boxed{
\text{модель}
\rightarrow
\text{алгоритм}
\rightarrow
\text{программа}
\rightarrow
\text{серия расчётов}
\rightarrow
\text{верификация/валидация}
\rightarrow
\text{анализ}}.
\]

\newpage
